% Related work and comparisons. Sources: HTML §7; refs/docs/model_arch.md §6
% (design comparison); AmphionASR 1.0 PDF §3.5, §5.4; the papers of MT-LLM, SpeakerLM, DiCoW,
% TS-SUPERB, M2MeT 2.0, NOTSOFAR-1, Personal VAD, TST-Bench and the two ISCA TSE papers
% (refs/notes/papers/INDEX.yaml, added 2026-11-01). The CocktailASR-1 protocol
% note (refs/docs/plan.md §8) lives in the experimental setup.
% Placed after the experiments, as in the HTML and refs/docs/plan.md §5.

\section{Related Work and Comparisons}
\label{sec:related}

\paragraph{Evaluation.}
Common evaluations for speech separation and TSE are synthesized from read speech: WSJ0-2mix~\cite{hershey2016deep}, Libri2Mix~\cite{cosentino2020librimix}, and WHAM!~\cite{wichern2019wham} and WHAMR!~\cite{maciejewski2020whamr}, which add noise and reverberation to WSJ0-2mix.
TSE usually enrolls another utterance of the same speaker and reports signal-to-distortion ratios on mixtures in which the target is present~\cite{zmolikova2019speakerbeam, zmolikova2023neural}.
Multi-talker and target-speaker recognition often use LibriSpeechMix~\cite{kanda2020serialized}, synthesized from LibriSpeech.
Real meeting corpora such as AMI~\cite{carletta2005ami}, AISHELL-4~\cite{fu2021aishell4}, AliMeeting~\cite{yu2022m2met} and CHiME-6~\cite{watanabe2020chime6} are mainly used for diarization and multi-talker transcription and do not themselves define enrollments or target-absent samples.
Neither do the speaker-attributed ASR benchmark of M2MeT~2.0~\cite{liang2023second} or the far-field office meetings of NOTSOFAR-1~\cite{vinnikov2024notsofar}.
TS-SUPERB~\cite{peng2025tssuperb} conditions TS-ASR on an embedding of the enrollment but evaluates it on Libri2Mix.
An absent target has mainly been studied outside recognition.
Personal VAD~\cite{ding2020personal} labels each frame as non-speech, target-speaker speech or non-target speech.
Speaker extraction has been trained and evaluated with the target present or absent in mixtures of one and two talkers~\cite{borsdorf2021universal, zhang2023speaker}, and Delcroix et al.~\cite{delcroix2022listen} studied false extractions in TSE when the target is silent.
TST-Bench~\cite{lee2026who} asks a system to tag enrolled speakers and reject unknown ones in long synthetic conversations.
In \tsbench, target-absent samples make up nearly half of the benchmark, and enrollments and absent samples are also defined on real meeting and call recordings.

\paragraph{Models.}
Early TS-ASR work added speaker conditioning to encoder--decoder recognizers such as Whisper~\cite{radford2023robust}: prompt tuning~\cite{ma2023extending}, learnable speaker queries that extract prompts from the enrollment~\cite{guo2024sqwhisper}, or frame-level diarization outputs that modulate the encoder~\cite{polok2024target}.
DiCoW~\cite{polok2024dicow} develops the last into masks of silence, target, non-target and overlapping speech and needs no enrollment.
Among SpeechLLMs, MT-LLM~\cite{meng2024large} follows multi-talker instructions, one of which places a 3\,s clip of the target talker and 3\,s of silence before the mixture.
SpeakerLM~\cite{yin2025speakerlm} jointly diarizes and transcribes a conversation, attributing speech to registered speakers through their speaker embeddings and leaving out registered speakers who do not speak.
AmphionASR~1.0~\cite{amphion2026amphionasr} encodes the enrollment and the mixture separately and feeds them to the language model in sequence.
Xiaomi-CocktailASR-1~\cite{zhang2026xiaomi} and its precursor~\cite{zhang2025thinking} concatenate reference speech with the mixture into one input to the audio encoder and support an empty output when the target is absent as well as chain-of-thought reasoning.
Like CocktailASR-1, AmphionASR~2.0 lets the enrollment interact with the mixture inside the audio encoder, but it joins the two segments after the convolutional downsampling rather than on the waveform (\Cref{sec:model}).
\Cref{tab:design} summarizes these systems.

\begin{table}[!tbp]
  \caption{Speaker-conditioned recognizers related to AmphionASR~2.0.}
  \label{tab:design}
  \centering
  \footnotesize
  \setlength{\tabcolsep}{4pt}
  \begin{tabularx}{\linewidth}{@{}>{\raggedright\arraybackslash}p{2.2cm}>{\raggedright\arraybackslash}p{2.25cm}>{\raggedright\arraybackslash}X>{\raggedright\arraybackslash}p{2.3cm}>{\raggedright\arraybackslash}p{2.6cm}@{}}
    \toprule
    System & Model & How the target is specified & Absent target & TS-ASR evaluation \\
    \midrule
    AmphionASR 2.0 & About 1.7\,B: AuT encoder + Qwen3-1.7B & Enrollment audio (first 5\,s; random 1--20\,s in training), joined to the mixture by a learnable \septok after the encoder's convolution & Empty output; scored on paired absent samples & \tsbench: Mandarin, synthetic and real \\
    \addlinespace[2pt]
    AmphionASR 1.0~\cite{amphion2026amphionasr} & About 1.7\,B, plus a 2.6\,M hotword-retrieval adapter & Enrollment audio (3--5\,s) and mixture pass through the same encoder and enter the language model as two ordered spans & Empty output; 164 samples in which nobody speaks & In-house synthetic set \\
    \addlinespace[2pt]
    CocktailASR-1 \cite{zhang2026xiaomi} & About 8.9\,B: Data2Vec2 encoder~\cite{baevski2023efficient} + Qwen3-8B & Reference speech (random 1--4\,s), 1\,s of silence and the mixture, concatenated on the waveform & Empty output; separate absent test sets & LibriMix, LibriSpeechMix, AMI, AliMeeting \\
    \addlinespace[2pt]
    MT-LLM~\cite{meng2024large} & Llama-2-7B-chat with LoRA; Whisper and WavLM encoders & A 3\,s clip of the target and 3\,s of silence before the mixture, with an instruction & Not discussed & LibriSpeechMix; English and German \\
    \addlinespace[2pt]
    SpeakerLM~\cite{yin2025speakerlm} & Qwen2.5-7B-Instruct; SenseVoice-large encoder & Registered speakers: an embedding averaged over 2--10\,s clips, with the speaker's name & Registered speakers who do not speak are left out & Speaker-attributed transcription of AliMeeting, AISHELL-4, AISHELL-5 \\
    \addlinespace[2pt]
    DiCoW~\cite{polok2024dicow} & Whisper large-v3-turbo & No enrollment: frame-level masks from a diarization system & Not discussed & AMI, NOTSOFAR-1, Libri2Mix, LibriCSS \\
    \bottomrule
  \end{tabularx}
\end{table}

\subsection{Comparison with AmphionASR 1.0}
\label{subsec:related_v1}

\paragraph{Target-speaker recognition.}
On its in-house synthetic test set, AmphionASR~1.0 reports a target-present WER of 13.02\%, a miss rate of 4.13\% and a 6.10\% false-alarm rate on 164 target-absent samples in which nobody speaks~\cite{amphion2026amphionasr}.
That test set differs from \tsbench, so the numbers are not directly comparable.
The closest part of \tsbench is the ``no one speaking'' section, where the false-alarm rate of AmphionASR~2.0 is 0.8\%; \Cref{subsec:exp_conditions} shows what an evaluation restricted to target-present and silent samples misses.

\paragraph{General recognition.}
The two reports share 10 public test sets (\Cref{tab:vs_v1}).
Because their Mandarin normalizations differ, the table compares each version with the same reference model, Qwen3-ASR-1.7B, as measured in each report.
In English, 1.0 matches or slightly beats Qwen3-ASR-1.7B on LibriSpeech and trails it by 0.6--0.9 points on GigaSpeech and Common Voice, while 2.0 trails slightly on LibriSpeech, by less on GigaSpeech (0.40) and by 0.81 on Common Voice.
In Mandarin, the gap of 2.0 to Qwen3-ASR-1.7B is smaller than that of 1.0 on KeSpeech and Common Voice, and both versions are slightly better than Qwen3-ASR-1.7B on AISHELL-2.
On the two WenetSpeech subsets the gap of 2.0 is larger (1.01 points on test\_net, against 0.13 for 1.0).

\begin{table}[!tbp]
  \caption{Comparison with AmphionASR~1.0 on shared public test sets (CER / WER, \%). Each version is compared with Qwen3-ASR-1.7B as measured in its own report; negative differences mean better than Qwen3-ASR-1.7B.}
  \label{tab:vs_v1}
  \centering
  \small
  \setlength{\tabcolsep}{5pt}
  \begin{tabular}{@{}lcccccc@{}}
    \toprule
    & \multicolumn{3}{c}{AmphionASR 1.0 report~\cite{amphion2026amphionasr}} & \multicolumn{3}{c}{This report} \\
    \cmidrule(lr){2-4}\cmidrule(l){5-7}
    Test set & 1.0 & Qwen3-ASR & $\Delta$ & 2.0 & Qwen3-ASR & $\Delta$ \\
    \midrule
    LibriSpeech $\cdot$ test-clean & 1.59 & 1.61 & $-$0.02 & 1.69 & 1.63 & +0.06 \\
    LibriSpeech $\cdot$ test-other & 3.35 & 3.39 & $-$0.04 & 3.65 & 3.37 & +0.27 \\
    GigaSpeech & 9.53 & 8.68 & +0.85 & 9.09 & 8.70 & +0.40 \\
    Common Voice $\cdot$ English & 7.87 & 7.24 & +0.63 & 8.17 & 7.36 & +0.81 \\
    AISHELL-1$^\dagger$ & 1.34 & 1.58 & $-$0.24 & 0.38 & 1.48 & $-$1.10 \\
    AISHELL-2 & 2.65 & 2.77 & $-$0.12 & 2.53 & 2.56 & $-$0.04 \\
    KeSpeech (all dialects pooled) & 5.71 & 5.21 & +0.50 & 5.19 & 4.97 & +0.22 \\
    Common Voice $\cdot$ Mandarin & 5.21 & 5.06 & +0.15 & 5.02 & 5.03 & $-$0.01 \\
    WenetSpeech $\cdot$ test\_net & 5.17 & 5.04 & +0.13 & 5.49 & 4.49 & +1.01 \\
    WenetSpeech $\cdot$ test\_meeting & 6.30 & 5.85 & +0.45 & 5.51 & 4.69 & +0.82 \\
    \bottomrule
  \end{tabular}
  \par\smallskip
  {\footnotesize Both reports normalize English with the Whisper English normalizer but normalize Mandarin differently: the 1.0 report removes punctuation, unifies full- and half-width characters and converts Arabic numerals to Chinese characters, and this report additionally removes filler words and unifies erhua. Absolute Mandarin numbers are therefore not comparable across the two reports. $^\dagger$AISHELL-1 shares text with AISHELL-2, which 2.0 is trained on; the 2.0 entry in this row is optimistic.\par}
\end{table}

\FloatBarrier
